% Generated draft: requires mathematical and source review.
\documentclass[11pt]{article}
\usepackage[margin=25mm]{geometry}
\usepackage{fontspec}
\setmainfont{Latin Modern Roman}
\usepackage{amsmath,amssymb,mathtools,unicode-math}
\setmathfont{Latin Modern Math}
\usepackage{enumitem,needspace,xurl,hyperref}
\setlength{\parindent}{0pt}
\setlength{\parskip}{5pt}
\setlength{\emergencystretch}{3em}
\newcommand{\N}{\mathbb{N}}
\newcommand{\Z}{\mathbb{Z}}
\newcommand{\Q}{\mathbb{Q}}
\newcommand{\R}{\mathbb{R}}
\newcommand{\C}{\mathbb{C}}
\newcommand{\F}{\mathbb{F}}
\newcommand{\A}{\mathbb{A}}
\newcommand{\PP}{\mathbb{P}}
\newcommand{\E}{\mathbb{E}}
\newcommand{\eps}{\varepsilon}
\newcommand{\dd}{\,\mathrm{d}}
\begin{document}
\section*{UnsolvedMath ID: 30000550}
\textbf{Source problem code:} catalogue entry
\subsection{Definitions and mathematical statement}
Let \(T^2=\R^2/\Z^2\). For a smooth Riemannian metric \(g\), a nonconstant closed geodesic is a periodic curve satisfying \(\nabla_{\dot\gamma}\dot\gamma=0\), considered without redundant changes of parameter. Free homotopy classes of oriented loops in \(T^2\) are indexed by \(a\in\Z^2\).

Under the usual geometric-distinctness reading, let \(\mathcal G_g(a)\) be the set of geometric images of nonconstant closed \(g\)-geodesics that have a representative in class \(a\). The source question becomes
\[
\exists g\text{ a Riemannian metric on }T^2\quad
\exists n_0\in\N,\ n_0\geq1\quad
\forall a\in\Z^2,\qquad |\mathcal G_g(a)|<n_0.
\]
The inequality is strict and the same finite bound applies to every class for one fixed metric. It concerns all closed geodesics in the class, not only length-minimizing ones, and no genericity or curvature condition is imposed.

The source does not define its counting convention. Counting parameterized maps would artificially create many copies by reparameterization, while counting geometric images, prime geodesics, or iterated geodesics up to a chosen equivalence can give different statements. Constant curves are excluded in the displayed conventional reading; otherwise the contractible class would contain every constant map. The wording ``every homotopy class'' includes \(a=0\) unless the source intends only nontrivial classes, a restriction not added here. Orientation reversal would identify \(a\) with \(-a\) under an unoriented convention. The appropriate convention must be confirmed before the metric-existence question is considered fully specified.

The supplied assessment checked 2026-08-21 reports the uniform-per-class bound unresolved in this formulation. Results giving many closed geodesics in total do not by themselves decide a uniform bound within each of infinitely many classes; the supplied related reference is not treated as resolving that distinction.
\subsection{Short English statement}
Can one metric on the two-torus bound, by one finite number, all closed geodesics within each free homotopy class? The source must fix geometric versus iterated counting and treatment of the contractible class; total infinitude of geodesics is a different issue.
\subsection{Sources}
\begin{itemize}
\item Local UnsolvedMath registry, permanent ID 30000550, catalogue entry, \emph{Uniform Bounds on Closed Geodesics in Torus Homotopy Classes}.
\item \emph{Geometric Group Theory, Hyperbolic Dynamics and Symplectic Geometry} (2007), Oberwolfach Reports, pp. 1991--2058, DOI \url{https://doi.org/10.4171/owr/2006/33}.
\item Supplied reference \emph{Closed geodesics and the first Betti number}: \url{https://arxiv.org/abs/2407.02995}.
\item Supplied literature triage checked 2026-08-21; no exact counting convention or resolving theorem is reproduced.
\end{itemize}
\end{document}
